\documentclass[openany, a4, 11pt]{book}

% Documentul va fi tradus in limba romana
\usepackage[romanian]{babel}

% Scriem cu diacritice... 
\usepackage[utf8x]{inputenc}

% Vom folosi si un appendix simplut
\usepackage[toc,page]{appendix}

% Se importa si pachetul pentru URL-uri (acestea vor fi colorate cu albastru)
\usepackage{hyperref}
\usepackage[pdftex]{graphicx}    
\usepackage{subfig}
\usepackage{multicol}
\usepackage{color}
\usepackage{indentfirst}
\usepackage{pifont}
\usepackage[margin=2.8cm, vmargin={1.4in, 1.4in}]{geometry}
% \usepackage[bindingoffset=0.2in, left=1.0in, right=1.0in]{geometry}

% Pachetele necesare pentru crearea si afisarea formulelor matematice
\usepackage{mathtools}

% Folosim un pachet pentru antet si subsol
\usepackage{fancyhdr}
\pagestyle{fancy}

% Pachetele necesare pentru listarea de cod sursa
\usepackage{caption}
\usepackage{listings}
\usepackage{xcolor}

% Pachetele necesare pentru crearea unor casute cu diverse mesaje
\usepackage{styles/boxes}

\usepackage{array}
\usepackage{booktabs}
\usepackage{multirow}

\usepackage{wrapfig}
\usepackage{graphicx}
\usepackage{subfig}

\definecolor{gray}{RGB}{235,235,235}
\definecolor{codegreen}{rgb}{0,0.6,0}
\definecolor{codegray}{rgb}{0.5,0.5,0.5}
\definecolor{codepurple}{rgb}{0.58,0,0.82}
\definecolor{backcolour}{rgb}{0.95,0.95,0.92}

\renewcommand*\lstlistingname{Cod sursă}

% Numele paginii de Aneze din cuprins este redenumit
\renewcommand{\appendixtocname}{Anexă}

% Legaturile web vor fi colorate cu albastru
\newcommand{\myhref}[3][blue]{\href{#2}{\color{#1}{#3}}}%

% Alte modificari de configurare
\renewcommand{\headrulewidth}{1pt}
\renewcommand{\footrulewidth}{1pt}
\setlength{\headheight}{12pt}
\setcounter{secnumdepth}{6}
% Adancimea cuprinsului...
\setcounter{tocdepth}{3}
% Seteaza identarea unui paragraf
\setlength{\parindent}{1cm}

% Pentru cuprins se scoate culoarea si bordarea legaturilor
\hypersetup{urlbordercolor={1 1 1}, linkbordercolor={1 1 1}, citebordercolor={1 1 1}}

% De aici incepe documentul nostru!
\begin{document}

\newgeometry{margin=3cm}
\begin{titlepage}

% Titlul va aparea centrat
\begin{center}
 
\begin{figure}[htb]
\center{
\includegraphics[width=0.4 \textwidth]{images/freescale_logo.png}
}
\end{figure}
\noindent{\rule{8cm}{0.4pt}}\\[0.3cm]
\textsc{\Large STEM Innovation Champ}\\[5.0cm]

{\Huge DOCUMENTAȚIE}\\[0.4cm]
\textsc{\Large Utilă pentru realizarea unui robot de tip\\ LINE FOLLOWER}\\[5.0cm]

% Informatii despre autor si despre indrumator
\begin{minipage}{0.5\textwidth}
\begin{flushleft}
\large{\textbf{Ediție concurs}}\\
\large{a 2-a (martie 2015)}
\end{flushleft}
\end{minipage}
~
\begin{minipage}{0.4\textwidth}
\begin{flushright}
\large{\textbf{Revizie document}} \\
2.0\\
\end{flushright}
\end{minipage}\\[3.0cm]

% Alte informatii aditionale?
\noindent{\rule{8cm}{0.4pt}}\\[0.3cm] 
{\large martie 2015}\\[1cm] 

\end{center}

% Trecem la urmatoarea pagina...
\end{titlepage}
\restoregeometry

% Cine a contribuit la acest proiect?
\newgeometry{margin=5cm}
\thispagestyle{empty}
\section*{\centering{Listă de schimbări}}

\restoregeometry

% Sectiunea "Abstract"
\newgeometry{margin=5cm}
\thispagestyle{empty}
\begin{multicols}{2}
[
\section*{Cuvânt înainte,}
]
\footnotesize {
Documentul de față își propune să vă ajute pe parcursul desfășurării proiectului \textit{"Line Follower"} organizat de \textit{Freescale România} în parteneriat cu \textit{Junior Achievement România}.

Dorim să vă informăm, de la bun început, de faptul că limbajul sau anumiți termeni ce apar în cadrul acestui document s-ar putea să nu vă fie familiari. Chiar dacă acest lucru se întâmplă, nu trebuie să vă panicați pentru că multe din aceste informații vă vor fi prezentate la sedințele tehnice pe care le veți avea cu mentorii. De fapt, prima ședință chiar asta îți propune: să aducă întreagă echipa la același nivel de cunoștințe și să aibă un vocabular adecvat. În celelalte trei ședințe ne vom ocupa de scris cod și optimizat soluția, dar începutul e cel mai important.

Vă recomandăm să citiți acest document înainte de a avea loc discuțiile cu mentorii. În acest fel, veți putea să vă faceți o idee despre ceea ce vom realiza, iar dacă ceva nu vă este clar, vă invităm să pregătiți o listă de întrebări pentru mentorii voștri.

Ce vei găsi în materialul de față? În primul rând, se vor regăsi explicații teoretice / fizice prin care se demonstrează modul de funcționare a unor module hardware și cum acestea pot fi programate folosind o unealtă ajutătoare. În al doilea rând, vei fi introdus în lumea sistemelor embedded, vei învăța cum poți controla un dispozitiv prin intermediul unui algoritm transpus într-un limbaj de programare, precum C.

}
\end{multicols}
\begin{flushright}\textit{Mult succes,\\Mentorii Freescale}\end{flushright}
\restoregeometry

% Se afiseaza si un cuprins avansat
\tableofcontents

% Pagina cu lista de figuri
\renewcommand\listfigurename{Listă de figuri}
\listoffigures

% Sfarsitul documentului
\end{document}